\Needspace{32\baselineskip}
\section{Explicit focus insertion and realization of typed search positions}
\label{sec:realization}

This section advances the operational obligation in the remaining-work list.
A term-vertex coordinate now determines an actual focussed hypernet, and on a
hole-free graph of type $*\to\epsilon$ it determines an actual state. The
construction proves all incidence and typing laws, rather than assuming a
realization interface. Contracting the inserted focus recovers the original
ordered incidence lists. Search preserves the term-vertex condition, so every
typed coordinate search step has two such state realizations.

This is the insertion direction of the bridge. It does not yet erase an
arbitrary pre-existing focussed state, establish inverse laws modulo graph
isomorphism, or identify the coordinate relation with the paper's contextual
interaction rules. The endpoint theorem below must not be read as a proof of
those still-open adequacy statements. Substitution, compute transitions, and
the concrete specimen interpretation also remain open.

\Needspace{16\baselineskip}
\subsection{Splitting a vertex}
\begin{definition}[Focus insertion carriers and maps]\label{term:focus-insertion}
Let $G$ be a hypernet and let $v$ be a term vertex that is not an output. Extend
the internal-vertex carrier from $V$ to $V+1$ and the edge carrier from $E$ to
$E+1$. Write $z$ for the fresh internal vertex, $\iota$ for the inclusion of old
ports, and $f$ for the fresh edge. Boundaries are unchanged. Define
\[
 S_v(w)=\begin{cases}z&w=v,\\\iota(w)&w\ne v,\end{cases}
 \qquad
 C_v(\iota(w))=w,\quad C_v(z)=v.
\]
Here $S_v$ redirects old source occurrences of $v$, while $C_v$ contracts the
new focus. Old internal vertices retain their wire labels and $z$ has label
term. The use of a singleton carrier gives explicit finite enumerations.
\end{definition}
\lean{Graph/FocusInsertion}{10}{10}{Linear.FocusInsertion.Vertex}
\lean{Graph/FocusInsertion}{12}{12}{Linear.FocusInsertion.host}
\lean{Graph/FocusInsertion}{14}{14}{Linear.FocusInsertion.fresh}
\lean{Graph/FocusInsertion}{16}{20}{Linear.FocusInsertion.collapse}
\lean{Graph/FocusInsertion}{30}{32}{Linear.FocusInsertion.split}
\lean{Graph/FocusInsertion}{63}{65}{Linear.FocusInsertion.internalLabel}

\begin{lemma}[Map properties]
The inclusion $\iota$ is injective and never reaches $z$. The contraction
satisfies $C_v\iota=\mathrm{id}$ and $C_vS_v=\mathrm{id}$; consequently $S_v$
is injective. No image of $S_v$ equals $\iota(v)$. Every new port is either an
old port's image or $z$. Both inclusion and redirection preserve wire labels.
\end{lemma}
\begin{proof}
Case analysis on internal, input, and output constructors proves the inclusion
claims; on internals use injectivity and disjointness of sum tags. For
$C_vS_v$, split on $w=v$: the fresh port contracts to $v$ in the equality case,
and otherwise contraction undoes inclusion. Applying $C_v$ to equal redirection
images proves injectivity. An image cannot equal $\iota(v)$: in the equality
case it is fresh; otherwise inclusion injectivity would force $w=v$.
The port classification uses the two internal sum cases and the boundary
constructors; every inhabitant of the fresh singleton carrier equals its sole
member. Inclusion preserves labels by cases. Redirection preserves labels
because only $v$ changes image, and both $v$ and $z$ have term label.
\end{proof}
\lean{Graph/FocusInsertion}{22}{22}{Linear.FocusInsertion.collapse_host}
\lean{Graph/FocusInsertion}{24}{24}{Linear.FocusInsertion.host_ne_fresh}
\lean{Graph/FocusInsertion}{26}{27}{Linear.FocusInsertion.host_injective}
\lean{Graph/FocusInsertion}{34}{38}{Linear.FocusInsertion.collapse_split}
\lean{Graph/FocusInsertion}{40}{42}{Linear.FocusInsertion.split_injective}
\lean{Graph/FocusInsertion}{44}{49}{Linear.FocusInsertion.split_ne_host}
\lean{Graph/FocusInsertion}{51}{61}{Linear.FocusInsertion.vertex_cases}
\lean{Graph/FocusInsertion}{67}{68}{Linear.FocusInsertion.host_label}
\lean{Graph/FocusInsertion}{70}{75}{Linear.FocusInsertion.split_label}

\begin{definition}[Inserted incidence lists]
For every old edge $e$, use $S_v(\mathsf{source}_G(e))$ for its source list and
$\iota(\mathsf{target}_G(e))$ for its target list, mapping entry by entry.
The fresh edge has source $[\iota(v)]$, target $[z]$, and the chosen focus mode
as label. Old edge labels, including whole boxed bodies, are unchanged.
\end{definition}
\lean{Graph/FocusInsertion}{77}{79}{Linear.FocusInsertion.source}
\lean{Graph/FocusInsertion}{81}{83}{Linear.FocusInsertion.target}

\begin{lemma}[Incidence uniqueness and coverage]
The inserted incidence lists satisfy source and target uniqueness. A port
occurs in a source list exactly when it is not an output, and in a target list
exactly when it is not an input.
\end{lemma}
\begin{proof}
For two old source edges, pull equality of mapped vertices back along the
injective $S_v$ and apply source uniqueness in $G$. An old source cannot meet
the fresh edge's source because $S_v$ avoids $\iota(v)$. There is only one
fresh edge. The target proof is identical with inclusion in place of
redirection and disjointness from $z$ in place of avoidance of $\iota(v)$.

For source coverage, all mapped old sources are non-outputs: a redirected one
is internal, and any other one retains its old boundary kind. The fresh edge's
source is a non-output by the hypothesis on $v$. Conversely, an old image
$\iota(w)$ with $w=v$ is sourced by the fresh edge. If $w\ne v$, use coverage
in $G$ and its retained source edge. The fresh port $z$ is sourced by the old
edge that sourced $v$, which exists because $v$ was not an output.
For target coverage, retained targets remain non-inputs and the fresh target
is internal. Conversely an old non-input image inherits a target edge from
$G$, and $z$ is targeted by the fresh focus edge.
\end{proof}
\lean{Graph/FocusInsertion}{85}{104}{Linear.FocusInsertion.source_unique}
\lean{Graph/FocusInsertion}{106}{125}{Linear.FocusInsertion.target_unique}
\lean{Graph/FocusInsertion}{127}{154}{Linear.FocusInsertion.source_coverage}
\lean{Graph/FocusInsertion}{156}{174}{Linear.FocusInsertion.target_coverage}

\begin{theorem}[Well-formed explicit focus insertion]
The construction gives a finite typed hypernet with the same external
interface as $G$.
\end{theorem}
\begin{proof}
Enumerate each extended carrier by the sum of the old enumeration and the
singleton enumeration. Old incidence lists retain their no-repetition property
because both maps are injective; the new lists are singletons. Use the proved
uniqueness and coverage lemmas. For typing, mapping a port list and then reading
labels is the same as reading the old labels, by the label-preservation laws.
The new source and target both have term label, exactly the focus signature.
These checks populate every field of the concrete net record.
\end{proof}
\lean{Graph/FocusInsertion}{177}{226}{Linear.FocusInsertion.insert}

\Needspace{16\baselineskip}
\subsection{Focus certification, contraction, and holes}
\begin{theorem}[Unique focus and contraction]
If $G$ is globally focus-free, the inserted graph is focussed in the sense of
\S6.2. Contracting the new focus restores each original source and target list.
\end{theorem}
\begin{proof}
Choose the fresh edge as focus. Every old label remains globally focus-free;
the singleton fresh-edge carrier has no second member. The focus endpoints are
$\iota(v)$ and $z$, which are distinct. Its target is internal, so the focus is
not exposed. Its incidence lists and mode are those specified by construction.
To recover source lists, compose their map by $S_v$ with the map by $C_v$ and
use $C_vS_v=\mathrm{id}$. For target lists use $C_v\iota=\mathrm{id}$.
These statements concern the retained edges; the focus itself is deleted when
interpreting contraction. They establish recovery for the constructed family,
not a general erasure operation on arbitrary focussed graphs.
\end{proof}
\lean{Graph/FocusInsertion}{229}{245}{Linear.FocusInsertion.focussed}
\lean{Graph/FocusInsertion}{248}{255}{Linear.FocusInsertion.collapse_sources}
\lean{Graph/FocusInsertion}{257}{264}{Linear.FocusInsertion.collapse_targets}

\begin{lemma}[Hole freedom and boundary observations]
If $G$ is hole-free at every depth, so is its focus insertion. The explicit
focus is initial exactly when its mode is search and $v$ is an input; it is
final exactly when its mode is backtrack and $v$ is an input.
\end{lemma}
\begin{proof}
Represent a hypothetical hole occurrence by an outer edge and a label-local
position. If the edge is retained, its unchanged label gives a hole occurrence
in $G$, a contradiction. If it is the fresh edge, its focus label is not a hole
and has no children. This covers zero and successor depth alphabets.
For boundary observations, the mode is unchanged and the source is
$\iota(v)$. Inclusion preserves input status, by cases on the port constructor.
\end{proof}
\lean{Graph/FocusInsertion}{267}{286}{Linear.FocusInsertion.holeFree}
\lean{Graph/FocusInsertion}{289}{293}{Linear.FocusInsertion.boundary_iff}

\Needspace{16\baselineskip}
\subsection{Search preserves term typing}
\begin{lemma}[Typed operation data]
An operation view's source has term label. Mapping its argument list through
vertex labels gives exactly the list of arity-many term labels. Consequently
every argument vertex has term label.
\end{lemma}
\begin{proof}
Rewrite the operation's source-typing certificate by its singleton source and
operation label; the singleton signature forces the source label to be term.
For arguments, rewrite the target-typing certificate using the argument/thunk
append decomposition. Take an argument-length prefix on both sides. The
arity certificate ensures that this is exactly the entire argument-label list
on one side and the term prefix of the signature on the other. Finally map
membership of an argument into this repeated-term list: every member equals
term.
\end{proof}
\lean{Graph/SearchRealization}{14}{18}{Linear.SearchOperation.source_term}
\lean{Graph/SearchRealization}{21}{35}{Linear.SearchOperation.arguments_typed}
\lean{Graph/SearchRealization}{37}{41}{Linear.SearchOperation.argument_term}

\begin{theorem}[Term typing preservation]
If a search step starts at a term vertex, its result is at a term vertex.
\end{theorem}
\begin{proof}
A variable trigger does not move, so inherits the premise. Operation entry
selects either the first argument or the source when there are no arguments.
Operation return selects either the next argument or the source on completion.
Apply the operation data lemmas in each case. Neither operation case needs the
premise on the old vertex; its result's type follows from graph typing alone.
\end{proof}
\lean{Graph/SearchRealization}{46}{60}{Linear.SearchStep.preserves_term}

\Needspace{16\baselineskip}
\subsection{Actual states from typed coordinates}
\begin{definition}[Typed position and realization]\label{term:typed-position}
For a graph $G:*\to\epsilon$, a typed position is a search coordinate whose
vertex is certified to have term label. Such a vertex cannot be an output,
since the output index type is empty. If $G$ is globally focus-free and
hole-free, realize a typed position by inserting its mode at its vertex and
packaging the resulting graph, unique-focus certificate, and hole-freedom
certificate into the concrete state record.
\end{definition}
\lean{Graph/SearchRealization}{63}{64}{Linear.TypedPosition}
\lean{Graph/SearchRealization}{70}{74}{Linear.TypedPosition.not_output}
\lean{Graph/SearchRealization}{77}{81}{Linear.TypedPosition.realize}

\begin{definition}[Initial attachment]\label{term:initial-attachment}
The initial typed position is search at the sole input, whose wire label is
term by the interface. Realizing it is an initial state: the constructed focus
has search mode and its source is that input. More generally, the realized
state's initial and final predicates coincide with their coordinate boundary
predicates, by the insertion boundary lemma.
\end{definition}
\lean{Graph/SearchRealization}{84}{85}{Linear.TypedPosition.initial}
\lean{Graph/SearchRealization}{88}{89}{Linear.TypedPosition.initial_realize}
\lean{Graph/SearchRealization}{92}{95}{Linear.TypedPosition.realize_boundary}
This discharges the graph-construction part of initial attachment. Proposition~38
still needs the concrete observation interface and specimen witnesses required
by \texttt{initial\_lift}.

\begin{definition}[Advancing a typed position]
Given a search step from a typed position, equip its result coordinate with
the term-typing proof supplied by preservation. This gives the next typed
position on the same underlying graph.
\end{definition}
\lean{Graph/SearchRealization}{98}{100}{Linear.TypedPosition.advance}

\begin{theorem}[Realization of search endpoints]
For every search step from a typed position on a focus-free, hole-free graph
$G:*\to\epsilon$, realize the initial position and its typed successor as
states $s,t$. Their explicit focus modes are the corresponding coordinate
modes. At each endpoint, collapsing the inserted focus recovers every
original edge's source and target list in $G$.
\end{theorem}
\begin{proof}
Use the realization constructor for the input and for the typed successor.
Their focus modes are correct by the definition of insertion. Apply the two
incidence-recovery theorems at each coordinate. The common underlying graph
is unchanged; only the position and mode of the inserted focus may differ.
The statement supplies actual state endpoints and recovery certificates, not
a new definition of a contextual graph-rewrite step between arbitrary states.
\end{proof}
\lean{Graph/SearchRealization}{104}{122}{Linear.TypedPosition.step_realizes}

\Needspace{22\baselineskip}
\Needspace{16\baselineskip}
\subsection{Boundary and interior regression checks}
The existing value corolla is focus-free because its operation label cannot
be a focus label. Realizing its initial typed position gives the explicit
state \texttt{initialValue}. Advancing by its passive nullary search step and
realizing the result gives \texttt{finalValue}. Their initial and final
certificates follow from their modes and common input source. The initial
state's distinct focus endpoints are guaranteed by its focussed certificate.
\lean{Graph/RealizationExamples}{8}{10}{Linear.Examples.value_focusFree}
\lean{Graph/RealizationExamples}{13}{14}{Linear.Examples.initialValue}
\lean{Graph/RealizationExamples}{17}{19}{Linear.Examples.finalValue}
\lean{Graph/RealizationExamples}{21}{22}{Linear.Examples.value_boundary_states}
\lean{Graph/RealizationExamples}{25}{27}{Linear.Examples.initial_value_has_distinct_endpoints}
The two-hole chain is also globally focus-free, because neither hole label can
be a focus. Insert search at its internal vertex: the first edge still targets
the old vertex, the second edge now starts at the fresh vertex, and the new
focus connects the old vertex to the fresh one. Unfolding the incidence maps
proves all four equalities. This focus is not initial because its source is
internal. Unlike the value example, this test deliberately retains holes and
therefore tests focussed hypernets rather than closed states.
\lean{Graph/RealizationExamples}{29}{31}{Linear.Examples.twoHoles_focusFree}
\lean{Graph/RealizationExamples}{35}{42}{Linear.Examples.interior_insertion}
\lean{Graph/RealizationExamples}{45}{49}{Linear.Examples.interior_not_initial}
